\section{Related Work}
\label{sec:related}

\paragraph{Anchoring, shortcuts, and memory traps.}
That imitation and VLA policies latch onto training-distribution regularities rather than
task-causal features is well documented. Causal confusion~\citep{dehaan2019causal} shows
behavioral cloning exploits nuisance correlates of expert actions; object-aware
regularization~\citep{park2021object} targets spurious object cues. In the VLA setting,
``memory traps''~\citep{afi2025memory} name the failure of reproducing memorized
trajectories instead of adapting to a moved scene, and propose an affordance-field fix;
proprioceptive state has been identified as a shortcut that harms spatial
generalization~\citep{proprio2025}. We take the \emph{existence} of anchoring as given.
Our contribution is not the phenomenon but an instrument that measures, on the
\emph{successful} rollouts, whether the reach is grounded or anchored---a question these
works do not pose.

\paragraph{Evaluating spatial generalization.}
Recent work argues that success-only benchmarks give incomplete or misleading pictures.
LIBERO-X~\citep{liberox2026} introduces a robustness litmus with progressive perturbations
and a capability decomposition; ``Robust Skills, Brittle Grounding''~\citep{robustskills2026}
perturbs object locations and decomposes success into skill-execution versus
instruction-grounding metrics. Spatial-tolerance and success-region analyses characterize
\emph{where} a policy works~\citep{mobipi2025,n2m2025,manibox2024}. These measure success
degradation, capability axes, or spatial extent. None measures, on the rollouts that
succeed, whether the endpoint moved with the displaced object or stayed at the canonical
location---the basin-versus-grounding distinction (\Cref{tab:delta}). Our readout is the
manipulation-reach analogue of interventional acceptance tests for
attribution~\citep{adebayo2018sanity}: it actively perturbs the cause and checks the
effect.

\paragraph{Counterfactual data and fine-tuning geometry.}
We use counterfactual-augmentation fine-tuning purely as a controlled intervention, in the
lineage of MimicGen~\citep{mandlekar2023mimicgen} and DemoGen~\citep{demogen2025}; we make
no data-generation claim. Our boundary analysis (\Cref{sec:boundary}) connects to work on
fine-tuning geometry---linear-probing versus full fine-tuning~\citep{kumar2022finetuning},
last-layer retraining for spurious correlations~\citep{kirichenko2023last}, and robust
weight interpolation~\citep{wortsman2022wise}---which studies classification, not embodied
reach.

\begin{table}[t]
\centering
\small
\caption{What related evaluation work measures, versus the reach readout. Only the readout
is conditioned on \emph{successful} rollouts and separates basin-tolerance from grounding.}
\label{tab:delta}
\begin{tabular}{lcccc}
\toprule
 & Perturbs & Measures reach & Conditioned on & Separates basin \\
 & object? & endpoint? & successes? & from grounding? \\
\midrule
Robust Skills, Brittle Grounding~\citep{robustskills2026} & yes & partial & no & no \\
LIBERO-X~\citep{liberox2026} & yes & no & no & no \\
AFI / Memory Traps~\citep{afi2025memory} & yes & no & no & no \\
\textbf{Reach readout (ours)} & \textbf{yes} & \textbf{yes} & \textbf{yes} & \textbf{yes} \\
\bottomrule
\end{tabular}
\end{table}
